\documentclass[11pt]{article}
\usepackage[a4paper,margin=28mm]{geometry}
\usepackage{amsmath,amssymb,amsthm,booktabs,microtype}
\usepackage[hidelinks]{hyperref}

\newtheorem{theorem}{Theorem}
\newtheorem{proposition}{Proposition}
\newtheorem{remark}{Remark}

\title{A Computer-Assisted All-Orders Certificate for a Six-Field
Flux Subsystem on a Calabi--Yau Threefold}
\author{Evgeniy Agafonov\\
\small Independent Researcher, Sochi, Russia\\
\small \href{mailto:my@eagafonov.ru}{my@eagafonov.ru}\\
\small ORCID: \href{https://orcid.org/0009-0005-9059-5291}{0009-0005-9059-5291}}
\date{Draft of 12 August 2026}

\begin{document}
\maketitle

\begin{abstract}
We develop a computer-assisted framework for proving existence and local
uniqueness of a supersymmetric flux-sector root on an explicit toric
Calabi--Yau threefold with $h^{2,1}=5$ and $h^{1,1}=81$.  The system contains
the axiodilaton and five complex-structure fields.  A real $12\times12$
Krawczyk operator certifies a root of a finite instanton system, while exact
Chow arithmetic and geometric-series majorants control the omitted GKZ
degrees through periods, F-terms, and their Jacobian.  The finite 557-degree
source sector is included in a recentered Arb evaluator with a box-wide
interval enclosure, and an independent verifier checks the complete degree
partition and the finite-plus-infinite Krawczyk inclusion. Five light quotient classes are identified
as persistent rigid rational complete-intersection curves, as a supporting
geometric result rather than part of the central proof.
\end{abstract}

\section{Context and contribution}

Algorithmic algebraic geometry and numerical homotopy continuation have long
been used to locate and enumerate flux vacua
\cite{Gray2006,Mehta2011,MartinezPedrera2013}. More recent work has made
large-modulus mirror symmetry computationally effective even for Calabi--Yau
threefolds with many moduli \cite{Demirtas2024}, while interval Krawczyk
operators provide practical certificates for isolated numerical solutions
\cite{Breiding2024}. The explicit geometry and flux choice used here descend
from the ancillary data of \cite{Demirtas2021}. Exact flux landscapes can
also arise on special symmetry loci \cite{Grimm2024}. Recent analyses stress
that instanton corrections can change the location and even the multiplicity
of flux solutions near large-complex-structure boundaries
\cite{Chauhan2026}.

The contribution of this paper is to join these ingredients into one
end-to-end certificate. Every degree is assigned either to an exact finite
Chow computation or to an analytic infinite-tail majorant; the resulting
period bounds are propagated through the mirror map, the six F-terms, their
real Jacobian, and the final Krawczyk operator. To our knowledge, the cited
flux-vacuum literature does not provide this combination of an explicit
multi-parameter compact Calabi--Yau model, all-orders GKZ accounting, and an
interval proof of local uniqueness. This is a deliberately narrower claim
than a proof of a complete physical vacuum.

\section{Claim and scope}

Let $F:\mathbb R^{12}\rightarrow\mathbb R^{12}$ be the real form of the six
complex equations $D_IW=0$, with $I$ running over the axiodilaton and five
complex-structure moduli.  The intended result is strictly local.

\begin{theorem}[Certified six-field local theorem]\label{thm:target}
For the flux and basis data specified in the archived artifact, the complete
closed-sector GKZ system has exactly one zero in the declared componentwise
box of radius $10^{-30}$ around the archived center.
\end{theorem}

The theorem does not claim global uniqueness, simultaneous stabilization of
the 81 K\"ahler moduli or open-string fields, uplifted stability, or a complete
physical vacuum.

% Generated by scripts/render_certificate_table.py; do not edit by hand.
% certificate sha256: 1205e43b3c447a54d08ca4804c6af8e0c6b541b55805316f533bee42b818911c
\begin{table}[t]
\centering
\small
\begin{tabular}{@{}lr@{}}
\toprule
Certified quantity & Value\\
\midrule
Complex fields / real equations & 6 / 12\\
Finite source rows & 557\\
Low-total representations & 8150\\
Partition $Z/S/M/U$ & 7150/527/21/452\\
Box radius $r$ & $1.000000\times10^{-30}$\\
$\lVert C F(x_0)\rVert_\infty$ & $3.727099\times10^{-69}$\\
$\lVert I-CJ(X)\rVert_\infty$ & $7.573429\times10^{-4}$\\
Finite Krawczyk image & $7.573429\times10^{-34}$\\
Finite strict margin & $9.992427\times10^{-31}$\\
$\lVert C_{\rm new}\rVert_\infty$ & $1.446862\times10^{23}$\\
Maximum all-orders addition & $1.081536\times10^{-53}$\\
Tail / finite margin & $1.082356\times10^{-23}$\\
Strict real components & 12\\
Authenticated provenance files & 16\\
\bottomrule
\end{tabular}
\caption{Machine-generated summary of the admitted certificate. The
displayed decimals are derived from exact rational bounds.}
\label{tab:certificate}
\end{table}


\section{Model and equations}

The compactification is the toric hypersurface dataset \texttt{5-81-3213} of
\cite{Demirtas2021}, with $(h^{2,1},h^{1,1})=(5,81)$. We use the primitive
integral divisor basis
\begin{equation}
 (J_1,J_2,J_3,J_4,J_5)=(D_2,D_3,D_4,D_5,D_9),
 \qquad \int_X c_2(X)\wedge J_a=(16,8,12,20,8)_a.
\end{equation}
The independent Chow reconstruction has full column rank over $\mathbb Q$;
the alternative compatible divisor tuple is related to the displayed one by
a unimodular matrix. In the displayed basis, the nonzero symmetric triple
intersections are
\begin{align}
&(\kappa_{111},\kappa_{112},\kappa_{114},\kappa_{122},
  \kappa_{124},\kappa_{134},\kappa_{144},\kappa_{222})
 =(-2,2,-2,-2,2,2,-2,2),\\
&(\kappa_{224},\kappa_{244},\kappa_{334},\kappa_{344},
  \kappa_{444},\kappa_{445},\kappa_{455},\kappa_{555})
 =(-2,2,-2,2,-4,2,-2,2).
\end{align}

Let $\Phi=(\tau,z^1,\ldots,z^5)$, where $\tau$ is the axiodilaton and the
$z^a$ are complex-structure coordinates. The integral flux vectors, in the
archived symplectic ordering, are
\begin{align}
 f&=(-2,4,-10,-2,3,-3,0,3,-5,2,-2,-5),\\
 h&=(0,-5,5,-4,-1,5,0,0,0,0,0,0).
\end{align}
Writing $f=(f_B,f_A)$ and $h=(h_B,h_A)$ in two six-component blocks and
$\Pi=(X^I,F_I)$ with $X^0=1$ and $X^a=z^a$, our convention is
\begin{equation}
 W=\sqrt{\frac{2}{\pi}}\left[(f_B-\tau h_B)^T X
 -(f_A-\tau h_A)^T F\right],\qquad \mathcal F_I=D_IW.
\end{equation}
The real system is ordered as
\begin{equation}
 x=(\operatorname{Re}\tau,\operatorname{Re}z^1,\ldots,
 \operatorname{Re}z^5,\operatorname{Im}\tau,
 \operatorname{Im}z^1,\ldots,\operatorname{Im}z^5).
\end{equation}
Its recentered witness begins
\begin{align}
x_0\simeq(&-12,-19.5,-29.5,-15,-15,-5;\nonumber\\
          &19.81991494,32.20736178,48.72395756,
          24.77489367,24.77489367,8.258297891),
\end{align}
with the full decimal center serialized in the certificate.

In the large-complex-structure chart the genus-zero prepotential is written
as
\begin{equation}
 \mathcal G(z)=-\frac16\kappa_{abc}z^az^bz^c
 +\frac12 A_{ab}z^az^b+\frac1{24}c_{2,a}z^a
 +\frac{\chi\,\zeta(3)}{2(2\pi i)^3}
 -\frac1{(2\pi i)^3}\sum_{\beta>0}n^{(0)}_\beta
   \operatorname{Li}_3(e^{2\pi i\beta\cdot z}),
\end{equation}
where $\chi=-152$, and $A_{ab}$ is the integral quadratic convention fixed
by the archived symplectic basis. The period blocks used above are
$X=(1,z^a)$ and $F=(2\mathcal G-z^a\mathcal G_a,\mathcal G_a)$.

\section{Finite-system interval certificate}

For a center $x_0$, interval box $X=x_0+[-r,r]^{12}$, and numerical inverse
witness $C$, define
\begin{equation}
 K(X)=x_0-CF(x_0)+(I-CJ(X))(X-x_0).
\end{equation}
The Arb calculation reports the conservative scalar bounds
\begin{align}
 r &=10^{-30}, &
 \lVert CF(x_0)\rVert_\infty &<3.73\times10^{-69},\\
 \lVert I-CJ(X)\rVert_\infty &<7.58\times10^{-4}.
\end{align}
Consequently the finite-system image radius is bounded by
\begin{equation}
 3.73\times10^{-69}+7.58\times10^{-4}r<r.
\end{equation}
This proves existence and local uniqueness for the implemented finite system.
The attachment also contains all twelve componentwise inclusions; the scalar
relaxation and the exact proof summary are reported in
Table~\ref{tab:certificate}.

\section{GKZ degree decomposition}

The secondary-cone monoid is decomposed into a zero-action direction and five
positive-action Hilbert generators.  Through transverse total eight, 8150
representations are classified as shown in Table~\ref{tab:partition}.

\begin{table}[ht]
\centering
\begin{tabular}{lr}
\toprule
Category & Representations\\
\midrule
Chow zero & 7150\\
Finite source table & 527\\
Displayed primitive multicover & 21\\
Uncovered nonzero Chow sector & 452\\
\midrule
Total & 8150\\
\bottomrule
\end{tabular}
\caption{Low-total degree ledger. The artifact contains the explicit disjoint
membership lists and their canonical digest.}\label{tab:partition}
\end{table}

The 452 representations reduce to 188 distinct nonzero Chow charges.  Exact
rational coefficient norms are used for this low sector.  All totals at least
nine enter a geometric majorant, whose action is bounded below by
$66.8922\ldots$.  The resulting uncovered-sector bound on a second logarithmic
period derivative is $1.743\times10^{-77}$.

\section{The finite-source-sector admission lemma}

\begin{proposition}[Certified finite-sector lemma]\label{lem:finite}
On a declared complex neighborhood containing the final Krawczyk box, the
difference between the complete 557-degree Chow-reduced period jet and the
finite baseline used in the interval solve has rigorous componentwise bounds
through derivative order two.  Those bounds are propagated through the
mirror map, symplectic periods, F-terms, and real Jacobian.
\end{proposition}

The full finite system shifts the original four-term center by approximately
$2.93\times10^{-20}$, so the old radius-$10^{-30}$ box is not retained. The
complete finite sector is instead included in the Arb evaluator and certified
in a new radius-$10^{-30}$ box around the recentered root. The inverse mirror
map is simultaneously certified by a five-component contraction argument.

\section{Certified all-orders perturbation}

For the already covered complement of the finite source sector, propagation
gives
\begin{align}
 \lVert\Delta F\rVert_\infty &<8.543\times10^{-77},\\
 \lVert\Delta J\rVert_\infty &<2.810\times10^{-73},\\
 \max_i \rho_i &<1.082\times10^{-53}.
\end{align}
Here the final verifier computes, using rational arithmetic,
\begin{equation}
 \rho_i=\sum_{j=1}^{12}|C_{{\rm new},ij}|
       \bigl(\Delta F_j+r\Delta J_j\bigr).
\end{equation}
All 144 upper bounds for $|C_{{\rm new},ij}|$ are serialized, and their exact
row sums give $\lVert C_{\rm new}\rVert_\infty<1.447\times10^{23}$.
Together with Proposition~\ref{lem:finite} and the exhaustive membership
ledger, these inequalities give the stated all-orders six-field theorem.

\begin{proof}[Proof of Theorem~\ref{thm:target}]
Let $X=x_0+[-r,r]^{12}$. The Arb witness encloses the complete finite
557-row system and its Jacobian on $X$, including the inverse mirror map, and
gives
\begin{equation}
 \lVert C F_{\rm fin}(x_0)\rVert_\infty
 +\lVert I-CJ_{\rm fin}(X)\rVert_\infty r
 \le 7.574\times10^{-34}.
\end{equation}
The explicit low-total ledger and the high-total geometric majorant exhaust
the complementary GKZ degrees. Their value and derivative bounds hold on the
radius-$0.1$ logarithmic domain, which contains $X$ after recentering. Exact
propagation through $C_{\rm new}$ adds at most
$1.082\times10^{-53}$ to each real component of the Krawczyk image. Hence
\begin{equation}
 K_{\rm full}(X)\subset\operatorname{int}X
\end{equation}
in all twelve components. The Krawczyk theorem therefore gives a zero of the
full closed-sector map in $X$ and uniqueness within the declared box.
\end{proof}

\section{Five supporting rigid curves}

Five minimum-volume quotient-basis classes have unique representatives among
the 3570 pairs of prime toric divisors:
\begin{equation}
 D_{15}D_{24}X,\quad D_{13}D_{84}X,\quad D_9D_{81}X,\quad
 D_{10}D_{80}X,\quad D_{13}D_{82}X.
\end{equation}
For each representative, exact intersection data and face combinatorics give
a primitive genus-zero class with normal bundle
$\mathcal O_{\mathbb P^1}(-1)\oplus\mathcal O_{\mathbb P^1}(-1)$.  The
restricted hypersurface equation is a primitive binomial whose coefficients
remain nonzero in the all-orders logarithmic coefficient box, so the curves
persist there. The interpretation of these classes follows the toric
complete-intersection analysis of \cite{Demirtas2024}.

This statement is intentionally not exhaustive.  The complete-intersection
audit contains 54 non-potent divisor-pair records representing 52 distinct
full classes.  The five curves support the geometric interpretation of the
selected quotient basis but do not prove an all-orders GV vanishing theorem.

\section{Reproducibility and trusted base}

The release contains a small Python verifier with no third-party dependency.
It recomputes strict scalar and componentwise inclusion inequalities, checks
scope constraints, authenticates upstream artifacts by SHA-256, and fails if
the finite-sector or degree-partition witnesses are absent.  The heavy
generators use Arb, exact rational Chow arithmetic, Normaliz, and CYTools.
Approximate roots and preconditioners are witnesses; interval inclusion is the
proof.

Negative tests alter preconditioner entries, tail bounds, one radius, one
category count, the local-uniqueness scope, and the geometric exhaustion
flag.  Each alteration must invalidate the certificate.  A separate
dependency-free verifier recomputes the twelve tail contributions and the
final inclusion inequality using exact rational arithmetic.

The table included in the manuscript is generated directly from
\texttt{data/certificate.json} and
\texttt{data/recentered\_tail\_propagation.json}. It is therefore a build
artifact rather than a second manually maintained statement of the numerical
claim.

The exact public release used for this article is version v1.0.2, archived
under the \href{https://doi.org/10.5281/zenodo.21906026}{Zenodo concept DOI}
(accessed 12 August 2026); the source is maintained in the
\href{https://github.com/megazord78-pixel/interval-flux-vacuum-5-81}
{public GitHub repository}.

\section{Discussion}

Computer-assisted proofs can bridge the gap between formal instanton
expansions and numerical flux vacua, but only when every series sector is
accounted for by a disjoint and exhaustive certificate.  The present workflow
provides a strict enclosure of the complete finite source sector, a sharp
majorant for the uncovered infinite tail, and an exact componentwise
propagation of that tail through the recentered preconditioner.  Their strict
combined inclusion proves the stated local six-field theorem.  No claim about
the full physical vacuum is made by this focused result.

\bibliographystyle{unsrt}
\bibliography{references}
\end{document}
